\chapter{Kết luận và hướng phát triển}
\label{chap:conclusion}

\section{Mức đáp ứng mục tiêu}
Đề tài đã mô hình hóa được một lượt gửi xe từ lúc nhận thẻ đến lúc thu tiền và trả thẻ. Kết quả chính nằm ở sự liên kết giữa những dữ liệu dùng chung: một thẻ chỉ thuộc một lượt đang gửi, phí gắn với thời gian của chính lượt đó, lịch sử không mất khi thẻ được dùng lại và xe chưa ra vẫn được khôi phục sau khi mở lại chương trình.

Năm nội dung trọng tâm của học phần đều có phần cài đặt cụ thể. Yêu cầu được chuyển thành điều kiện trước và sau thao tác. Các lớp xe, lượt gửi, bãi và thời gian giữ những trách nhiệm khác nhau. Kế thừa và đa hình tập trung ở quy tắc phí; nhóm ngoại lệ phân biệt các nguyên nhân thất bại. Những đoạn mã trong báo cáo được lấy từ chương trình để việc giải thích có thể đối chiếu trực tiếp khi bảo vệ bài.

Kết quả 45 ca C++ và sáu kịch bản console hỗ trợ các kết luận trong phạm vi đã thử. Đặc biệt, các ca hủy trả, báo giá cũ và lỗi ghi tệp kiểm tra được thời điểm dữ liệu được phép thay đổi. Chúng bổ sung cho kiểm thử mức tiền, vì một mức phí đúng vẫn chưa đủ nếu lượt bị kết thúc nhầm hoặc doanh thu được ghi trước khi xác nhận. Luồng mất thẻ bổ sung một đường tra cứu theo biển số, lưu thông tin sinh viên/chủ xe và thu phí thủ công; các ca T36--T45 xác nhận lượt vẫn đóng đúng và dữ liệu đọc lại đầy đủ.

\section{Những lựa chọn cần cân nhắc khi mở rộng}
Cách lưu bằng tệp văn bản giúp nhóm dễ theo dõi giao dịch và tạo lại các tình huống lỗi, nhưng toàn bộ lịch sử phải được duyệt và ghi lại sau mỗi thay đổi. Nếu dùng với lượng dữ liệu lớn hơn, cần thay lớp lưu trữ và cách tra cứu trước khi nghĩ đến việc đổi giao diện.

Quy tắc qua đêm cũng cần được thống nhất trước khi tự động hóa. Đề bài mới chốt ba mức phí trong cùng ngày, nên chương trình yêu cầu nhân viên nhập phí và lý do cho lượt qua ngày. Một công thức tự động chỉ nên được thêm khi chính sách đã rõ.

\section{Hướng phát triển}
Trước mắt, nhóm cần thống nhất quy tắc qua đêm và cách xử lý khi lưu thất bại. Sau đó có thể bổ sung kho thẻ, tài khoản nhân viên và nhật ký thao tác. Khi quy mô tăng, lớp lưu trữ có thể chuyển sang cơ sở dữ liệu và hỗ trợ nhiều cổng; phần cứng như đầu đọc thẻ hay camera chỉ nên thêm sau khi quy trình nghiệp vụ đã ổn định.

Kiến trúc tách giao diện, dịch vụ và lưu trữ tạo điểm bắt đầu thuận tiện cho các thay đổi này, nhưng mỗi bước vẫn phải được kiểm thử lại. Trong phạm vi học phần, sản phẩm hiện tại đủ để trình diễn một quy trình có dữ liệu, có điều kiện lỗi và có các cơ chế OOP có thể giải thích rõ ràng.
